\section{Conclusion and Discussion}\label{sec:conclusion}
Our research has demonstrated the power of non-prehensile manipulation and whole-body interaction in enabling robotic manipulators to operate effectively despite locked multi-joint failures. By leveraging these two key strategies, we have shown that our approach can significantly increase the manipulable area, allowing the robot to maintain up to 92\% of the nominal reachable area, even when the gripper is unusable. Furthermore, our sim-in-the-loop planner effectively utilizes kinodynamic maps to complete manipulation tasks during LMJ failures. 

We believe future work can: i) explore and improve the completeness of our motion planner \cite{kleinbort2018probabilistic}; ii) expand the portfolio of motion primitives, such as pushing and rolling; iii) extend our approach to more precise manipulation tasks and other robot embodiments, like mobile manipulators; and iv) improve the simulation physics to match the real world more closely. Furthermore, we believe that applying optimization-based or learned approaches to joint failures can be fruitful areas of future research. This research paves the way for robots to operate independently for extended periods, even in the face of significant failures.
 